\section{One field, integer certificates and the all-degree bound}\label{sec:counts}

\begin{theorem}[Finite grid certificate]\label{th:finite}
Let $2^r\mid N$, $2^s\mid M$, $r,s\ge1$ and $N+(1+p)M\le D+1$. Write $R_{kl}^{(p)}$ for \eqref{eq:pure} or \eqref{eq:carrier} at $S=D+1$. There exists one field \eqref{eq:field} having at least
\begin{equation}\label{eq:gridcount}
 L_p(D;N,M,r,s)=\sum_{k=1}^r\sum_{l=1}^s
       \frac{NM}{2^{k+l}}\big(R_{kl}^{(p)}-1\big)_+
\end{equation}
distinct hyperbolic limit cycles, where $x_+=\max(x,0)$.
\end{theorem}
\begin{proof}
Fix $\lambda=\pi$ and, for $p=1$, fix one $\kappa$ outside the finite union in Theorem~\ref{th:rank}. Consider a diagonal block of rank $d\ge2$. Its germs at $h=0$ have a basis with distinct orders $n_1<\cdots<n_d$, obtained by elimination. The Wronskian has leading coefficient proportional to $\prod_{i<j}(n_j-n_i)$ and is nonzero at a sufficiently small \emph{positive} point $h_0$. There the map from coefficients to the Taylor polynomial through degree $d-1$ is invertible. Prescribe the Taylor polynomial $\prod_{i=1}^{d-1}(h-h_0-\delta i)$ for $\delta>0$. The chosen coefficients stay bounded as $\delta\to0$. Rescaling $h=h_0+\delta u$ and dividing by $\delta^{d-1}$ gives $\prod_{i=1}^{d-1}(u-i)+O_{C^1}(\delta)$ for bounded $u$. This yields $d-1$ simple positive zeros, on disjoint compact intervals strictly inside the annulus. Lift the linear combination to one form $\omega_{kl}\in\W_{kl}(S)$. The common cell sign in \eqref{eq:angles} gives the same zeros on every cell of the layer. Set $\omega_{kl}=0$ if $d\le1$.

Now take $\omega=\sum_{k,l}\zeta^{k+l}\omega_{kl}$. In a cell of $J_k\times I_l$ only indices $k'\ge k$, $l'\ge l$ can contribute. After division by $\zeta^{k+l}$, the first period is the diagonal period plus $O_{C^1}(\zeta)$ on its compact intervals. With one sufficiently small $\zeta>0$ all the selected zeros persist, since there are finitely many cells and intervals. Fix this one real polynomial form. On each interval \eqref{eq:period} divided by $-\mu$ now has a simple zero $h_\mu=h_*+O(\mu)$ by the implicit function theorem. The multiplier $1-\mu L_C'[\omega](h_*)+O(\mu^2)$ is positive and different from one for small $\mu>0$. Choose disjoint compact collars around the selected ovals and take the minimum of their finitely many smallness thresholds. All these cycles belong to the same $X_\mu$ of degree at most $D$. The layer sizes give \eqref{eq:gridcount}.
\end{proof}

For a square grid $N=M=mq$, where $q=2^r$, take depth $r$ in both variables (also when $m$ is even). The following closed forms exhibit the degree cost of the carrier.
\begin{corollary}\label{cor:symmetric}
For integers $r,m\ge1$, the unweighted and weighted square grids respectively give
\begin{align}
 \HH(2mq-1)&\ge m^2\left(\tfrac34 q^2 r^2-\tfrac12 q^2+\tfrac32 q-1\right),\label{eq:sym0}\\
 \HH(3mq-1)&\ge m^2\left(\tfrac{15}{8}q^2 r^2-\tfrac14 q^2+\tfrac54 q-1\right).
 \label{eq:sym1}
\end{align}
\end{corollary}
\begin{proof}
For a nonnegative integer $j$,
\[
 P_{kl}(jN)=j^2 2^{k+l-3}+\mathbf1_{k=l}\,j2^{k-2}.
\]
Indeed the odd positive indices $u,v$ obey $u/2^k+v/2^l\le j$. For $k=l$ put $h=j2^{k-1}$ and get $h(h+1)/2$. For $k<l$, each odd $u<j2^k$ leaves $j2^{l-1}-2^{l-k-1}u$ choices for $v$. Summing over $u$ uses the fact that the first $h$ odd integers sum to $h^2$. Interchange $k,l$ for the other half. Equality along the diagonal cutoff accounts for the second term.

In \eqref{eq:pure} the combination is $2P(2N)-2P(N)$, giving $3\cdot2^{k+l-2}+\mathbf1_{k=l}2^{k-1}$. In \eqref{eq:carrier} it is $2P(3N)-3P(N)$, giving $15\cdot2^{k+l-3}+\mathbf1_{k=l}3\cdot2^{k-2}$. All ranks exceed one. Multiplying each by $N^2/2^{k+l}$ makes the main contribution constant across the $r^2$ blocks. In \eqref{eq:gridcount} the diagonal corrections sum to $N^2(1-1/q)/2$ and $3N^2(1-1/q)/4$, respectively, and subtracting one per cell costs $N^2(1-1/q)^2$. This yields the two formulas.
\end{proof}

In these two grids the leading contributions per $(D+1)^2r^2$ are $(3/4)/2^2=3/16$ and $(15/8)/3^2=5/24$. Thus the additional rank pays for the extra $M$ degrees of the carrier even after both counts are expressed in the final degree.

\begin{proof}[Proof of Theorem~\ref{th:main}]
For $S=D+1\ge16$, let $b=1+\lfloor\log_2 S\rfloor$ be its binary length. Choose $q$ as the largest power of two with $bq\le S$, and put $m=\lfloor S/(3q)\rfloor$. These integer operations give an explicit witness at every such budget. Since $S/(2b)<q\le S/b$, we have $m\ge1$, $r=\log_2 q\ge1$, $3mq-1\le D$, and
\[
 N=mq=S/3+O(S/\ln S),\qquad
 r=(\ln S-\ln\ln S)/\ln2+O(1).
\]
The right side of \eqref{eq:sym1} is $(15/8)N^2 r^2+O(S^2)$. In $N^2$ the rounding error $O(S^2/\ln S)$ costs $O(S^2\ln S)$ after multiplication. Squaring $r$ gives $[(\ln S)^2-2\ln S\ln\ln S]/(\ln 2)^2+O(\ln S)$, since $(\ln\ln S)^2=O(\ln S)$. Substitution, $(15/8)/9=5/24$ and $S=D+1$ give \eqref{eq:main}. The witness uses the budget $3mq-1$ even when the prescribed degree is between grid degrees.
\end{proof}

\subsection{Comparison at the same degree}
For $q=3^r$, $r\ge2$, Qin and Sun's Theorem~1.1 produces $4r^2q^2/9-(q-1)^2$ hyperbolic cycles at degree at most $4q-5$. The Hamiltonian has degree $2q$ and the perturbation sets the final cost. Substitution of the final degree gives coefficient $1/[36(\ln3)^2]$ along the subsequence. Their Corollary~1.2 uses monotonicity between successive seeds, giving $1/[324(\ln3)^2]$ for every $D\ge31$, with a loss of order $D^2\ln D$ relative to that leading term~\cite[Theorem 1.1, Corollary 1.2 and Section 6]{qin}.

The Chebyshev replication of Eshkobilov, Kadyrov and Mamayusupov copies $L$ cycles of degree $d$ to $m^2L$ cycles at degree at most $m(d+1)-1$~\cite[Theorem 1]{replication}. Each branch is a local orbital equivalence, so the multiplier is preserved or inverted. Applied to the ternary seeds, this recovers the subsequence coefficient for all sufficiently large $D$, with the same logarithmic rounding term as in \eqref{eq:main}. Lu's Proposition~5.2 gives this combination explicitly. The mechanical theorem gives the larger coefficient $1/[8(\ln2)^2]$ with loss $O(D^2\ln D)$. Thus the comparisons involve both the degree quantifier and the precision below the leading term~\cite[Theorem 1.1 and Proposition 5.2]{mechanical}.

For a finite comparison, let $G(D)$ be the maximum of \eqref{eq:gridcount} over the two carrier choices and admissible even $N,M$, using their full binary depths. Shorter depths only omit nonnegative summands. The four comparison families in Table~\ref{tab:finite} are those specified in Lu's finite-degree supplement~\cite[Sections 1--3]{supplement}. Its common copy and unused-degree rule takes a seed ledger $b$ to $\clos[b](D)=\max_{d\le D}b(d)\lfloor(D+1)/(d+1)\rfloor^2$. The classical column uses its four-root Han--Li tree and listed direct seeds, the two-source column starts empty, the ternary column uses the seeds above, and the mechanical column starts from the binary-prefix renewal of~\cite[Proposition 3.3]{mechanical}. The classical ledger records bounds for $H$ as stated there. The other columns count hyperbolic cycles.

\begin{table}[tb]
\caption{Finite certificates at total degree at most $D$. The triple is one maximizing $(p,N,M)$ for the grid. A dash means that the ternary ledger has no seed at that degree.}\label{tab:finite}
\centering\small
\begin{tabular}{r r l r r r r}
\toprule
 $D$&$G(D)$&$(p,N,M)$&Classical&Two-source&Ternary&Mechanical\\
\midrule
3&3&$(0,2,2)$&13&0&---&0\\
5&11&$(0,2,4)$&37&0&---&4\\
7&45&$(0,4,4)$&74&0&---&20\\
11&140&$(0,4,8)$&153&80&---&72\\
23&1,161&$(0,8,16)$&833&440&---&676\\
31&2,967&$(0,16,16)$&1,801&1,320&80&1,940\\
63&18,735&$(0,32,32)$&9,065&7,912&320&12,756\\
95&48,071&$(0,32,64)$&21,705&20,648&720&30,676\\
191&275,535&$(1,64,64)$&104,681&114,152&2,880&176,724\\
319&875,631&$(1,64,128)$&280,073&394,472&40,256&652,116\\
\bottomrule
\end{tabular}
\end{table}

At $D=5$, the symmetric weighted certificate gives eight cycles and the rectangular unweighted one gives eleven. The mechanical recurrence has half-degree three, prefixes $1,3$ and one birth contribution $4(3-1-1)=4$. The published reversible-center construction already gives $H(5)\ge37$~\cite{reversible}. At $D=11$ the classical ledger likewise exceeds the grid. On the finite range $5\le D\le350$, direct enumeration gives $G(D)>\clos[b_{\rm mech}](D)$ at each of the 346 degrees. These comparisons refer to the specified ledgers.

The verification files enumerate these integer bounds, reconstruct Taylor matrices at rational sample parameters and check the frequency collisions in Lemma~\ref{lem:shell}. The proof of Theorem~\ref{th:rank} determines the dimension for arbitrary admissible sizes. Its scope is the two displayed carriers and the stated blocks. Optimality among general positive polynomial multipliers and a further logarithmic factor lie outside these conclusions.
